\section{Full POMDP Specification}
\label{app:pomdp}

This appendix specifies the POMDP components used by the focal agent.

\subsection{Hidden state factors and observations}

As described in Section~\ref{sec:partner_models}, the focal agent's generative model for each partner $k$ maintains three hidden factors:
\begin{align}
s_k^{\text{type}} &\in \{\text{cooperator, reciprocator,} \nonumber\\
  &\phantom{s_k^{\text{type}} \in \{}\;\text{exploiter, random}\} && \text{(hidden behavioral type)} \label{eq:s_type}\\
s_k^{\text{stance}} &\in \{\text{trusting, neutral, hostile}\} && \text{(hidden stance)} \label{eq:s_stance}\\
s^{\mathrm{own}} &\in \mathcal{A}=\{0,1,2,3,4,5\} && \text{(executed investment level)}
\end{align}
The own-investment state is \emph{deterministic bookkeeping}: the agent always knows its executed investment level, and this factor makes that knowledge available to the observation model. It is not a psychological hidden variable. The joint hidden state per partner spans \(4 \times 3 \times 6 = 72\) configurations.

Observations are partner response and payoff:
\begin{align}
o_k^{\text{act}} &\in \{\text{cooperate, defect}\} & &\text{(partner response)}\\
o_k^{\text{pay}} &\in \{5,6,7,8,9,10,10.5,11,11.5,12,12.5\} & &\text{(focal agent's realized payoff)}
\end{align}

\subsection{Partner-response likelihood}

The core of the generative model is the \emph{partner cooperation likelihood table}, which defines the probability the partner cooperates given type and stance:
\begin{table}[H]
\centering
\caption{Cooperation likelihood $P(\text{cooperate} \mid s^{\text{type}}, s^{\text{stance}})$.}
\label{tab:coop_likelihood}
\begin{tabular}{lccc}
\toprule
Partner type & Trusting & Neutral & Hostile \\
\midrule
Cooperator & 0.95 & 0.80 & 0.55 \\
Reciprocator & 0.90 & 0.70 & 0.30 \\
Exploiter & 0.70 & 0.35 & 0.10 \\
Random & 0.60 & 0.50 & 0.35 \\
\bottomrule
\end{tabular}
\end{table}

This table is the foundation of the likelihood modality $A^{\text{act}}_k$. Entries define
\begin{equation}
P\!\left(o_k^{\mathrm{act}}=\mathrm{cooperate}\mid s_k^{\mathrm{type}}, s_k^{\mathrm{stance}}\right),
\label{eq:a_act_coop}
\end{equation}
with $P(o_k^{\mathrm{act}}=\mathrm{defect}\mid\cdot)=1-P(o_k^{\mathrm{act}}=\mathrm{cooperate}\mid\cdot)$. The modality is uniform over $s^{\mathrm{own}}$: partner response depends only on inferred type and stance. This operationalizes what the agent considers a predictable partner---a cooperative reciprocator in a trusting stance has high predicted cooperation; an exploiter in a hostile stance has low predicted cooperation.

\subsection{Payoff likelihood}

Payoff enters as a second observation modality \(A^{\mathrm{pay}}_k = P(o^{\mathrm{pay}}_k \mid s^{\mathrm{own}},s^{\mathrm{type}}_k,s^{\mathrm{stance}}_k)\). For investment \(a=s^{\mathrm{own}}\), endowment \(E=10\), and multiplier \(M=3\), the deterministic payoff map is
\begin{equation}
\pi(a,o_k^{\mathrm{act}})=
\begin{cases}
E-a+\frac{Ma}{2}, & o_k^{\mathrm{act}}=\mathrm{cooperate},\\
E-a, & o_k^{\mathrm{act}}=\mathrm{defect}.
\end{cases}
\label{eq:graded_payoff_map}
\end{equation}
Because the partner response is not directly controlled by the focal agent, the payoff likelihood marginalizes over the partner-response likelihood:
\begin{align}
&P(o^{\mathrm{pay}}_k=v \mid s^{\mathrm{own}},s^{\mathrm{type}}_k,s^{\mathrm{stance}}_k) \nonumber\\
&\quad =
\sum_{o^{\mathrm{act}}\in\{\mathrm{cooperate},\mathrm{defect}\}}
\mathbb{I}\!\left[\pi(s^{\mathrm{own}},o^{\mathrm{act}})=v\right]
P(o^{\mathrm{act}}_k=o^{\mathrm{act}} \mid s^{\mathrm{type}}_k,s^{\mathrm{stance}}_k).
\label{eq:a_pay_marginal}
\end{align}
This makes higher investment beneficial when cooperation is likely and costly when defection is likely, while keeping payoff preferences inside the generative model.

\subsection{Transition dynamics}

\paragraph{Type drift ($B_k^{\mathrm{type}}$).}
Partner type is uncontrollable and drifts slowly:
\begin{equation}
P(s_k^{\mathrm{type}\prime}=j\mid s_k^{\mathrm{type}}=i)=
\begin{cases}
1-p_{\mathrm{switch}} & j=i,\\
p_{\mathrm{switch}}/(K_{\mathrm{type}}-1) & j\neq i,
\end{cases}
\label{eq:b_type}
\end{equation}
with $K_{\mathrm{type}}=4$ and default $p_{\mathrm{switch}}=0.05$. Scheduled betrayal scenarios set $p_{\mathrm{switch}}=0$.

\paragraph{Stance dynamics ($B_k^{\mathrm{stance}}$).}
Stance transitions depend on the focal agent's graded investment $a$ by interpolating between a high-investment/trust-building endpoint and a low-investment/withholding endpoint:
\begin{equation}
\rho(a)=\frac{a}{|\mathcal{A}|-1}=\frac{a}{5},
\qquad
B^{\mathrm{stance}}(a)=\rho(a)B^{\mathrm{high}}+(1-\rho(a))B^{\mathrm{low}}.
\label{eq:b_stance}
\end{equation}
High investment shifts dynamics toward trust building; low investment or withholding shifts dynamics toward trust damage; recovery from hostility is slower than collapse.

\begin{table}[H]
\centering
\caption{Endpoint stance transition dynamics. The focal agent's graded investment interpolates between $B^{\mathrm{high}}$ and $B^{\mathrm{low}}$ according to Eq.~\eqref{eq:b_stance}.}
\label{tab:stance_transition}
\begin{tabular}{lccc}
\toprule
\multicolumn{4}{c}{\textbf{High investment / trust-building endpoint }$B^{\mathrm{high}}$} \\
\midrule
To $\setminus$ From & Trusting & Neutral & Hostile \\
\midrule
Trusting & 0.90 & 0.30 & 0.05 \\
Neutral & 0.10 & 0.60 & 0.35 \\
Hostile & 0.00 & 0.10 & 0.60 \\
\midrule
\multicolumn{4}{c}{\textbf{Low investment / withholding endpoint }$B^{\mathrm{low}}$} \\
\midrule
To $\setminus$ From & Trusting & Neutral & Hostile \\
\midrule
Trusting & 0.10 & 0.05 & 0.02 \\
Neutral & 0.50 & 0.35 & 0.18 \\
Hostile & 0.40 & 0.60 & 0.80 \\
\bottomrule
\end{tabular}
\end{table}

These dynamics instantiate the asymmetry of trust: sustained high investment gradually builds trust, while sustained withholding rapidly damages trust; recovery from hostility is slow.

\paragraph{Own-investment bookkeeping ($B^{\mathrm{own}}$).}
The own-investment factor updates deterministically from the executed investment level:
\begin{equation}
P(s^{\mathrm{own}\prime}=a^\prime \mid s^{\mathrm{own}}, a)=\mathbb{I}[a^\prime=a].
\label{eq:b_own}
\end{equation}
This factor is not a psychological hidden variable; it makes the agent's executed investment level available to the payoff likelihood in standard factorized form.

\subsection{Preferences, priors, and policy priors}

\paragraph{Observation preferences ($C_k$).}
There is no direct preference over partner responses:
\begin{equation}
C_k^{\mathrm{act}} = [0,\,0] \quad \text{over } \{\mathrm{cooperate},\mathrm{defect}\}.
\label{eq:c_act}
\end{equation}
Payoff preferences ascend over the graded payoff support generated by Eq.~\eqref{eq:graded_payoff_map}:
\begin{equation}
C_k^{\mathrm{pay}} = \mathrm{log\_softmax}\!\left(\frac{[5,6,7,8,9,10,10.5,11,11.5,12,12.5]}{\tau}\right),
\label{eq:c_pay}
\end{equation}
where $\tau$ is a temperature parameter controlling preference sharpness. Instrumental motivation therefore enters through payoff observations and multi-step planning rather than direct preferences over partner responses.

\paragraph{Initial-state priors ($D_k$).}
Type, stance, and own-investment priors are
\begin{align}
D_k^{\mathrm{type}} &= \left[\tfrac{1}{4},\tfrac{1}{4},\tfrac{1}{4},\tfrac{1}{4}\right], \label{eq:d_type}\\
D_k^{\mathrm{stance}} &= [0.2,\,0.6,\,0.2], \label{eq:d_stance}\\
D^{\mathrm{own}} &= \mathrm{Unif}(\{0,1,2,3,4,5\}). \label{eq:d_own}
\end{align}

\paragraph{Policy priors ($E_k$).}
Policy priors are uniform over admissible graded investment sequences at the configured planning horizon. In partner-choice settings, partner selection is handled outside the per-partner POMDP by evaluating partner--investment policy branches with precision-adjusted policy scores before executing the chosen investment.

\subsection{Partner implementation}

Partner implementation follows Section~\ref{sec:trust_game}. Four environment-side partners maintain hidden type and stance and sample cooperate-or-defect responses from the cooperation likelihood in Table~\ref{tab:coop_likelihood}, undergoing investment-dependent stance transitions per Table~\ref{tab:stance_transition}. Partners do not perform variational inference or maintain affective precision estimates, isolating the focal agent's precision mechanism before extending to reciprocal active-inference partners.
